%%
%% Supplementary Information working draft
%% Chinese methods-first version aligned with the v0.8.0 formula boundary.
%% National heterogeneous results remain pending rerun.
%%

\documentclass[preprint]{elsarticle}

\usepackage[UTF8,fontset=fandol]{ctex}
\usepackage{amsmath,amssymb,bm}
\usepackage{booktabs}
\usepackage{array}
\usepackage{tabularx}
\usepackage{longtable}
\usepackage{xcolor}
\usepackage{url}
\usepackage{hyperref}
\usepackage[margin=2.3cm]{geometry}

\hypersetup{colorlinks=true,linkcolor=blue,urlcolor=blue}
\renewcommand{\arraystretch}{1.18}
\setlength{\LTpre}{4pt}
\setlength{\LTpost}{4pt}
\renewcommand{\thetable}{S\arabic{table}}
\renewcommand{\thefigure}{S\arabic{figure}}
\renewcommand{\theequation}{S\arabic{equation}}

\newcommand{\implemented}{\textcolor{green!40!black}{\textbf{[当前实现]}}}
\newcommand{\literature}{\textcolor{blue!70!black}{\textbf{[文献依据]}}}
\newcommand{\adapted}{\textcolor{violet!75!black}{\textbf{[文献适配]}}}
\newcommand{\newwork}{\textcolor{orange!80!black}{\textbf{[本研究新增--需作者判断]}}}
\newcommand{\pending}{\textcolor{red!70!black}{\textbf{[待补充/待验证]}}}
\newcommand{\boundary}{\textcolor{black!65}{\textbf{[研究边界]}}}
\newcommand{\Aset}{\mathcal A}
\newcommand{\Hset}{\mathcal H}
\newcommand{\Nset}{\mathcal N}

\journal{Supplementary Information working draft}

\begin{document}

\begin{frontmatter}

\title{工厂级与行业级分布式大语言模型价值分析：补充方法、公式与研究边界}

\author{研究工作稿}

\begin{abstract}
本补充材料给出工厂级、集团级和行业级大语言模型部署的统一价值分析方法。比较固定有效AI服务，而不固定GPU数量或年用电；在共同任务质量、时延、可靠性与数据边界下，联合表示任务调度、CPU/GPU路由、服务器装机与在线状态、设施功率、企业净购电、分布式能源和新增接入容量。当前企业成本目标包括服务器、到户电量费、最大需量费、分布式能源及模型生命周期成本，不包含电网升级成本。电网侧只报告AI情景相对匹配无AI反事实的新增接入容量，以及相对参考架构避免的容量和避免比例；不将MW折算为货币投资。本文同时给出制造业AI需求的自下而上估计、31行业异质性处理、工业负荷映射、计算与设施参数以及各项企业成本参数的估计方法和来源，并区分观测值、文献代理和作者情景假设。全国31行业乘以三种架构的联合CPU/GPU结果尚未重跑完成，因此本版SI只冻结方法与参数边界，不报告或解释旧口径全国数值。
\end{abstract}

\end{frontmatter}

\tableofcontents
\clearpage

\section{阅读说明与证据标注}

本SI不是代码接口说明，而是供审稿人复核研究设计、公式逻辑和推断边界的补充方法文件。代码仅是实现和复现载体，不决定章节结构。全文使用以下标注：

\begin{itemize}
  \item \literature：定义、公式或机制已有公开研究或标准支持；
  \item \adapted：基本机制来自文献，但为分布式LLM架构比较进行了改写；
  \item \implemented：已经进入当前主流程的约束、目标或输出定义；
  \item \newwork：本研究提出的账户、指标或组合，需要作者判断是否作为方法贡献；
  \item \pending：缺少数据、校准、全国重跑或外部验证，不能作为已完成结论；
  \item \boundary：用于防止把筛查结果解释成已实现工程、社会总成本或现实因果关系。
\end{itemize}

\boundary 当前活动方法版本为\texttt{v0.8.0}。只有C36--IF样例完成了联合CPU/GPU物理核查；全国31行业$\times$三架构结果及其下游图表尚未按该口径重跑。因此，历史全国结果不进入本版SI的结果表，也不能用下列新公式重新解释。

\section{研究问题、功能单位与比较边界}

这一节首先回答“各种部署架构究竟在比较什么”。工厂侧、集团池和行业单节点只有在完成同样AI服务时才可以公平比较；否则，较低的电力或成本可能只是来自少做了任务。因此，本文先固定有效服务、质量和SLA，再让硬件数量、执行时间、用电和成本随架构变化。

\subsection{研究问题}

本研究考察：在完成相同制造业AI服务的条件下，将计算部署在工厂、集团池或行业单节点，会如何改变计算资源、设施用电、企业购电峰值、新增接入容量和年度增量成本？该问题同时包含两个尺度：

\begin{enumerate}
  \item \textbf{工厂级价值}：AI负荷能否利用生产负荷的非同时性、任务柔性、屋顶光伏或储能，降低企业新增接入容量与最大需量；
  \item \textbf{行业级价值}：跨工厂池化能否减少重复装机和备用，同时避免形成过大的单点连接需求。
\end{enumerate}

\subsection{功能单位与服务等价}

\implemented 公平比较的功能单位是“在共同质量、时延、可靠性、隐私和数据边界下完成的一单位有效AI服务”。令$D_{ikt}$为工厂$i$的任务$k$在时段$t$到达的有效服务需求，$d^{a}_{iknt}$为架构$a$在计算节点$n$实际执行的服务量。对刚性任务，执行时刻等于到达时刻；对柔性任务，服务可以在允许窗口内移动，但总量必须守恒：

\begin{equation}
 \sum_n d^{a}_{iknt}=D_{ikt},\qquad k\in\mathcal K^{\mathrm{rigid}},
 \label{eq:rigid-service}
\end{equation}

\begin{equation}
 \sum_{n,t}d^{a}_{iknt}=\sum_tD_{ikt},\qquad k\in\mathcal K^{\mathrm{flex}},
 \label{eq:flex-service}
\end{equation}

并要求所有架构达到相同服务质量与SLA门槛。未完成服务在当前主比较中为零，而不是通过较低服务质量换取成本下降。

任务还需满足数据和控制边界。令$M_{ikn}\in\{0,1\}$表示工厂$i$的任务$k$是否允许在节点$n$执行，则

\begin{equation}
 0\leq d^a_{iknt}\leq M_{ikn}D_{ikt}.
 \label{eq:task-admissibility}
\end{equation}

当生产控制、敏感检索或数据驻留规则禁止任务离开工厂或集团边界时，对应的$M_{ikn}=0$。这种做法把硬约束表示为可行性，而不用一个低置的惩罚价格“买下”隐私或安全例外。\boundary 当前数据边界仍是按任务类型设定的结构情景，尚未由企业数据分级和真实合规规则校准。

\literature 以服务输出而非设备铭牌作功能单位，符合计算系统能效与AI生命周期研究中“单位有效工作”比较的基本原则。\adapted 本研究把该原则扩展到制造业任务、企业负荷和接入容量的联合比较。

\subsection{外部用电情景的作用}

制造业AI用电的8、14和28~TWh/年情景只用于检查服务需求和服务--计算转换参数的数量级：

\begin{equation}
 E^{\mathrm{AI,derived}}\in[8,28]\ \mathrm{TWh/yr}.
 \label{eq:external-envelope}
\end{equation}

\boundary 式~\eqref{eq:external-envelope}不是服务量定义，也不是要求不同架构具有相同年用电。固定服务后，PUE、冷备、在线空闲、硬件路由、装机粒度和调度差异都可以使各架构年用电不同。

若某些边缘任务在所有架构中始终保留在工厂，总AI用电可分为

\begin{equation}
 E^{\mathrm{AI,total}}_a
 =E^{\mathrm{edge,common}}+E^{\mathrm{choice}}_a.
 \label{eq:common-edge-partition}
\end{equation}

共同边缘部分在架构差值中可以抵消，但仍属于制造业AI总用电。\boundary 只有在设备、服务质量和运行方式也完全相同时，该项才能从架构差异中取消；当前尚未单独报告全国共同边缘设备存量。

\section{参数估计原则与证据等级}

这一节要处理的问题是：模型同时使用统计观测、厂商规格、跨行业代理和研究情景，如果不区分证据强度，读者就无法判断结果是由观测驱动，还是由假设驱动。本研究因此按“观测优先、机制分解、代理透明、情景检验”的顺序估计参数，并为每个关键输入保留来源和适用边界。参数不因已经进入模型就自动成为观测值。为便于作者和审稿人判断，证据分为四级：A级为官方统计、正式标准或直接适用的监管价格；B级为厂商规格、企业披露或较匹配的实测数据；C级为跨地区、跨行业或跨设备代理；D级为缺少直接证据的作者情景假设。由多个来源综合构造的参数仍按最弱关键环节评级。

令参数$\theta$的低、中、高情景分别为$\theta^{L}$、$\theta^{B}$和$\theta^{H}$。若参数有直接观测，基准值优先取与研究地区、年份和技术边界相符的观测；若只有区间证据，基准值取区间内工程上可解释的中值而不是机械平均；若只有机制方向，则保留为结构情景并以结果敏感性判断其重要性。所有金额统一到人民币和共同价格年；仅企业购置的服务器、光伏和储能等资本项按后文资本回收因子年化。

\boundary A--D等级表示当前研究中的可追溯性和适用性，不是对来源本身质量的评价。D级参数可以用于筛查，但不能作为经验事实陈述；若其改变架构排序，结论必须降级为条件性结论。

\section{制造业AI服务需求的自下而上估计}

这一节把“制造业需要多少AI”拆成可解释的业务活动。本文想捕捉的不只是企业是否声称使用AI，而是不同行业有多少合格任务、多少企业采用某类任务、多大比例的工作流被覆盖，以及每个事件需要多少有效模型服务。为此，需求先按行业、企业规模、地区和任务自下而上构建，再汇总到全国。

\subsection{需求分解与防止重复计算}

需求分解的核心是把“企业外延”、“任务覆盖”和“单次服务强度”分开。这样可以避免把同一采用效果既算入企业数，又算入工作流覆盖率。\adapted 对行业$i$、企业规模$s$、地区$r$、任务$k$和年份$y$，年度有效AI服务需求写为

\begin{equation}
 D_{isrky}=B_{isrk}
 A^{\mathrm{any}}_{isry}\theta_{isky}W_{isky}
 f_{isky}n_{isky}S_{isky},
 \label{eq:bottom-up-demand}
\end{equation}

其中$B_{isrk}$是AI出现之前就存在的合格业务驱动量；$A^{\mathrm{any}}$是企业使用至少一种AI的比例；$\theta$是已使用AI的企业中采用任务$k$的条件比例；$W$是该任务采用者内部被AI覆盖的工作流或事件比例；$f$是每个驱动量每年的业务事件频率；$n$是每个被覆盖事件形成的有效模型任务或调用数；$S$是每次调用在共同质量和SLA下的有效服务规模。全国任务量为

\begin{equation}
 D_{ky}=\sum_{i,s,r}D_{isrky}.
 \label{eq:national-bottom-up-demand}
\end{equation}

式~\eqref{eq:bottom-up-demand}的直觉是：先数出可能产生AI任务的业务对象，再逐层缩小到真正采用且被AI覆盖的事件，最后换算为有效服务。硬件性能、量化、推测解码和PUE不进入该式，而在服务--计算和计算--电力层处理，以免把需求变化和技术效率变化混为一件事。IBM对全球2,900名高管的调查仅用于约束$A^{\mathrm{any}}\theta W$的数量级，而不直接给出制造业服务量或2030预测\cite{ibm2025}。

\subsection{采用率和企业规模}

采用率层用来限定哪些企业有可能产生AI任务，而不直接决定任务量。“使用任何AI”是企业外延边界，具体任务采用是其中的条件份额；将二者分开可以保留行业和企业规模差异，同时防止六类任务的采用率失去解释。

\literature 2023年任何AI采用的行业外延边际来自第五次全国经济普查综合卷表4--9；规模以上制造业合计96,311家报告使用AI，对应460,054家规模以上制造企业的20.93\%。行业观测从C20的约11.4\%到C16的约36.6\%，其中C36、C39和C40分别约为29.1\%、32.8\%和33.2\%\cite{nbs2023}。该指标包括机器视觉、自然语言处理和机器学习等，不是生成式AI采用率，也不表示生产部署深度。因此

\begin{equation}
 A^{\mathrm{task}}_{isrky}
 =A^{\mathrm{any}}_{isry}\theta_{isky},
 \label{eq:task-adoption}
\end{equation}

而不为六类任务分别设彼此无约束的总体采用率。

企业规模结构采用经济普查表2--22和2--23中的大型、中型、小型和微型企业数量与从业人员。规模以上企业的行业边际为观测值；规模以下企业不直接继承20.93\%，而设置更宽的采用区间。地区边际可由表4--10校准。若行业、规模和地区边际不能同时直接观测，则以迭代比例拟合使汇总严格回到已知边际。

\newwork 2030采用率、条件任务份额和工作流覆盖率是从2023观测基准出发的情景，不是经济普查的外推预测。政策普及目标只作为上界合理性检查，不能直接赋给企业或任务。

\subsection{任务驱动量与行业差异}

各任务使用不同的物理或业务驱动量，不能统一用营业收入或token替代。表~\ref{tab:task-drivers}给出优先估计量；缺少直接数据时，才使用从业人员、企业数或产线数代理，并保留代理区间。

\begin{table}[htbp]
 \centering
 \caption{自下而上任务需求的优先驱动量。}
 \label{tab:task-drivers}
 \begin{tabularx}{\textwidth}{p{3.0cm}X p{5.1cm}}
  \toprule
  任务 & 优先驱动量$B$与事件频率$f$ & 主要行业差异来源 \\
  \midrule
  办公知识服务 & 知识工作者或活跃用户数$\times$使用日/年 & 从业人员结构、管理与研发人员占比 \\
  业务流程Agent & 订单、采购、销售或供应链工作流数$\times$事件频率 & SKU、供应链层级和跨系统流程复杂度 \\
  中央VLM异常复核 & 检测工位或摄像头$\times$异常事件率 & 只计算被基础视觉筛出的异常，不计算全部视频帧 \\
  预测维护与诊断 & 关键设备数$\times$诊断窗口频率 & 设备密度、停机损失和连续生产程度 \\
  排程与工艺优化 & 工厂/产线数$\times$重排或优化频率 & 批次、换线、扰动和工艺耦合程度 \\
  研发仿真与数字孪生 & 研发团队、工程工作站或仿真作业数$\times$作业频率 & 研发强度、产品复杂度及设计/运行型孪生占比 \\
  \bottomrule
 \end{tabularx}
\end{table}

31个制造业大类保留各自的任何AI采用率、规模结构、活动量、代表集团份额和负荷曲线，同时按生产机制归入六类原型，用于任务组合和缺失曲线回退。该分类只共享机制先验，不把行业数据合并为一个平均行业。

\begin{longtable}{p{3.4cm}p{4.0cm}p{4.3cm}p{3.0cm}}
 \caption{制造业行业原型、行业代码和重点任务。}\label{tab:industry-archetypes}\\
 \toprule
 原型 & 行业代码 & 重点任务 & 柔性/负荷处理 \\
 \midrule
 \endfirsthead
 \toprule
 原型 & 行业代码 & 重点任务 & 柔性/负荷处理 \\
 \midrule
 \endhead
 食品与冷链 & C13--C15 & 供应链Agent、质量复核、排程、冷链维护 & 班次与冷链负荷；质量闭环偏刚性 \\
 批次流程 & C16、C27、C29、C42 & 批次排程、合规知识、维护、仿真 & 批次间可移时，生产中任务偏刚性 \\
 轻工制造 & C17--C24、C41（不含已另列行业） & 视觉复核、订单Agent、排程 & 班次明显；中小企业权重较高 \\
 连续高耗能 & C22、C25、C26、C28、C30--C32 & 维护、工艺优化、运行型孪生 & 基础负荷较平；关键控制任务增加刚性 \\
 离散装备 & C33--C38、C43 & 视觉、排程、维护、研发仿真 & 产线/产品差异大，设计仿真柔性较高 \\
 电子洁净制造 & C39--C40 & 高频视觉、良率诊断、排程、研发仿真 & 高停机损失；视觉与排程增加刚性 \\
 \bottomrule
\end{longtable}

\pending 表~\ref{tab:industry-archetypes}是机制分组，不是统计分类的替代。C22同时具有轻工产品属性和连续高耗能过程，在负荷和任务参数中按连续过程处理；具体行业仍使用独立观测值。

\subsection{从当前基准到2030情景}

未来需求同时受采用企业增加、工作流覆盖加深、事件频率变化和单次任务复杂度上升影响。这些驱动往往相关，所以本节先展示机械乘积，再通过联合情景收缩极端组合。

对于任务$k$，当前到2030的机械增长可写为

\begin{equation}
 g_k^{\mathrm{raw}}
 =g^A_kg^W_kg^f_kg^n_kg^S_k.
 \label{eq:raw-growth}
\end{equation}

采用率、覆盖率、事件频率、Agent深度和上下文规模通常相关，不能把各自高值独立相乘。本研究对机械乘积作联合相关性收缩，形成表~\ref{tab:service-growth}的低、中、高情景。它们保持任务间相对强弱，并以8/14/28~TWh为软边界检查，但不是文献点预测。

\begin{table}[htbp]
 \centering
 \caption{当前采用的2030有效服务量联合增长情景。}
 \label{tab:service-growth}
 \begin{tabular}{lrrr}
  \toprule
  任务 & 低 & 基准 & 高 \\
  \midrule
  办公知识服务 & 12 & 20 & 40 \\
  业务流程Agent & 75 & 130 & 260 \\
  中央VLM异常复核 & 70 & 120 & 240 \\
  预测维护与诊断 & 52 & 90 & 180 \\
  排程与工艺优化 & 145 & 250 & 500 \\
  研发仿真与数字孪生 & 128 & 220 & 440 \\
  \bottomrule
 \end{tabular}
\end{table}

\newwork 表~\ref{tab:service-growth}全部为证据约束的联合情景参数（D级），不是六项独立事实。其作用是把采用、覆盖和单任务复杂度显式拆开并允许替换；若企业日志可提供任何一项，应直接替换对应因子而非整体缩放总服务量。

\subsection{代表集团和工厂数量}

代表集团和工厂数用于把行业需求放缩到一个可求解的节点，再恢复到行业层。模型想捕捉的是代表性集团占行业多大份额、该份额分布在多少工厂，而不是复刻某一真实企业的网络。

典型集团行业份额$s_i^{\mathrm{group}}$优先由同口径集中度$CR_N$估计：

\begin{equation}
 s_i^{\mathrm{group}}\approx\frac{CR_{N,i}}{N},
 \qquad M_i^{\mathrm{group}}=\frac{1}{s_i^{\mathrm{group}}}.
 \label{eq:group-share-estimation}
\end{equation}

集团工厂数$F_i$来自代表性龙头企业年报、官网披露的国内主要生产基地或工厂数；若行业缺少集中度，则用相邻细分行业、产量或营业收入份额形成低、中、高区间。这里的$s_i^{\mathrm{group}}$和$F_i$是用于构造典型节点的代表值，不是对某个真实集团的识别。证据较强的行业（例如汽车、钢铁）可采用行业集中度；分散行业和跨细分行业代理应按C/D级处理。

\section{工业负荷、任务时序与计算参数的估计}

有效服务量还不能直接决定电力结果：同一年度需求在一天中何时到达、能够延迟多久、与何种生产负荷叠加，以及由CPU还是GPU完成，都会改变峰值、装机和用电。这一节因此把年度服务转换为具有行业负荷背景、任务到达时间和硬件需求的逐时输入。

\subsection{工业基础负荷曲线}

基础负荷的作用是捕捉AI功率与现有生产用电的同时性。年度用电量决定规模，连续代表周决定班次、周末和峰谷形状；二者分开可以避免用少量用户曲线推断全行业电量。行业$i$的代表工厂负荷由归一化曲线$\ell_{it}$与年度用电量$E_i^{L}$结合：

\begin{equation}
 L_{it}=E_i^{L}\frac{\ell_{it}}{\sum_{t\in\mathcal T}\ell_{it}\Delta t}\frac{1}{w^{\mathrm{ann}}}.
 \label{eq:base-load-scaling}
\end{equation}

曲线选择遵循“同一中国行业实测曲线、专属外部行业曲线、相近生产原型”的证据阶梯。26个行业使用中国南方EWELD数据集中相近ISIC行业的15分钟实测用户曲线；C31使用韩国和UCI钢铁工厂曲线，C37使用德国运输设备行业曲线；C16、C25和C42因缺少同类样本而回退到六类生产原型\cite{eweld2023,uciSteel,ffeLoad}. 每个行业选择连续168小时代表周并聚合至小时，保留测量时序而非把工作日平均复制七次。

\boundary EWELD只含386个工业和商业用户，严格高完整度样本更少；外部钢铁和运输设备曲线还有跨国差异。因此这些曲线主要识别负荷形状和AI负荷非同时性，不用于推断中国行业年度总电量。年度电量$E_i^L$与时序形状的来源和不确定性应分开报告。

\subsection{任务到达、柔性和行业修正}

任务日内到达形状依据业务机制构造：办公和Agent跟随工作时段，质量复核和维护更多跟随生产活动，排程集中于班前、换线或扰动时刻，仿真和离线任务可形成后台批次。中央AI任务分为刚性、日内和批处理三类，份额满足

\begin{equation}
 \rho_k^{\mathrm{rigid}}+\rho_k^{\mathrm{intra}}
 +\rho_k^{\mathrm{batch}}=1.
 \label{eq:flexibility-shares}
\end{equation}

基准情景中，办公、Agent、中央VLM、维护、排程和仿真的刚性份额分别为0.70、0.55、0.75、0.35、0.25和0.10；日内截止窗为4--12小时，批处理截止窗为72小时。普通闭环机器视觉属于边缘任务并保持100\%刚性，不与中央VLM异常复核混合。

\newwork 上述比例是按实时控制、人工交互和后台批处理边界构造的结构先验（D级），不是制造业全国观测。连续流程、电子和高停机损失行业把视觉、维护与排程约10个百分点从柔性份额移至刚性份额；研发密集行业根据设计验证或运行型孪生的占比调整仿真柔性。该修正只改变执行时间，不能同时再提高任务总量。

\subsection{有效服务到CPU/GPU计算}

有效服务转换参数$a^{\mathrm{eff}}$由六项相关的推理效率证据作加权几何综合，得到高效率、基准和保守情景0.2549、0.3333和0.4169参考L20加速器小时/有效服务单位。来源覆盖生成式AI功率轨迹、推理服务与能耗基准，但并非目标双L20服务器的直接制造业实测。因此当前主情景取保守值，三值均应作敏感性而不能称为硬件标定结果。

任务硬件份额$q_{kh}$和相对时间$m_{kh}$按“相同模型/任务、共同质量门槛、P95/P99时延和成功吞吐”估计。MLPerf Inference和AWS Inference Recommender支持这一比较口径，并表明CPU、专用推理芯片和GPU都应通过负载测试选型，而不能按TDP或理论FLOPS直接换算\cite{mlperf,AWSrecommender}。办公与Agent以GPU大模型和CPU工具链混合，中央视觉取决于吞吐和时延，维护允许较高CPU/边缘份额，排程为CPU/GPU混合，仿真以GPU或加速器为主。

\pending 当前没有全国任务级CPU/GPU路由观测，故$q_{kh}$和$m_{kh}$仍为结构情景。正式结论至少应报告全GPU、证据约束混合路由和CPU高占比三种边界；目标硬件实测完成前，不能把路由节省解释为已实现采购节省。

\section{部署架构与代表尺度}

本节建模的是AI需求被集中到多大的计算节点，而不是把制造业本身的生产负荷重新安置到一处。希望捕捉的两种相反机制是：集中化可以池化任务和服务器，但也会把原本分散的AI功率聚合成更大的单点负荷。模型通过代表工厂、代表集团和行业单节点的缩放关系，在保持行业服务总量一致的同时改变空间集中程度。

\subsection{当前比较架构}

\implemented 主文企业内部部署比较的架构集合为

\begin{equation}
 a\in\Aset=\{\mathrm{IF},\mathrm{IG}_{1\mathrm{host}},\mathrm{IG}_{\mathrm{multisite}}\}.
 \label{eq:architectures}
\end{equation}

\begin{table}[htbp]
 \centering
 \caption{当前部署架构及其解释边界。}
 \label{tab:architectures}
 \begin{tabularx}{\textwidth}{p{2.0cm}p{3.1cm}X p{3.6cm}}
  \toprule
  架构 & 名称 & 物理含义 & 解释边界 \\
  \midrule
  IF & 工厂侧分布式 & 每个代表工厂分别配置CPU/GPU服务器，并与本厂基础负荷和DER共同运行。 & 工厂独立装机的分散边界；不额外计入专网和治理成本。 \\
  IG-1host & 集团单节点算力池 & 同一集团的工厂需求在一个代表成员工厂池化。 & 主文核心代表场景；表示集团集中配置，不表示全行业共享。 \\
  IG-multisite & 集团多节点协调 & 集团在多个成员工厂安装算力，可迁移任务联合选择执行时刻和工厂。 & 跨节点灵活性扩展；受数据、网络和时延边界约束。 \\
  \bottomrule
 \end{tabularx}
\end{table}

原II\_1host行业单节点退出Figure~2的企业内部三方案比较，仅保留为大型集中上界或云端资源足迹反事实。集团多节点协调与“同一企业同时使用自建设备和云端预留容量”的采购混合不能混称：前者改变集团内部计算节点网络，后者是在给定部署架构内分配本地与云端服务。

\boundary 当前Figure~2结构情景只对IF启用逐厂GPU/CPU整数装机；IG-1host与IG-multisite采用连续等效装机，三者逐时在线容量均连续。因此IF与IG-1host的差额同时反映逐厂离散采购和集团池化，不能全部归因于组织集中；IG-multisite与IG-1host采用相同定容口径，二者差额才用于识别跨节点调度价值。对只有一个直接EWELD用户的行业，成员工厂曲线由同一用户的不同完整周按星期--小时对齐构造，并明确标记为多周代理而非真实跨厂观测。

为识别AI与原生产负荷错峰的价值，只对核心IG-1host配对运行保留固定承载工厂负荷的`actual-load`和将该负荷设为零的`zero-load`。两个边界保持AI服务量、释放时刻、截止时间、CPU/GPU路由、价格和连续等效定容规则不变，并各自重新优化执行时间。定义

\begin{equation}
 V^{\mathrm{align}}_{\mathrm{IG-1host}}=C_{\mathrm{IG-1host}}^{\mathrm{zero,reopt}}-C_{\mathrm{IG-1host}}^{\mathrm{actual,reopt}}.
\end{equation}

$V^{\mathrm{align}}_{\mathrm{IG-1host}}>0$表示在当前模型边界下，固定承载工厂的既有生产负荷形状降低了AI增量成本。该值分解为服务器、AI电量、最大需量和新增接入容量差额。IF和IG-multisite不运行零负荷；跨节点价值另由实际负荷下IG-multisite相对IG-1host的差额识别。\boundary `zero-load`是AI独立负荷反事实，不是停产工厂场景；因两组均重新优化，该指标包含允许的调度和定容响应，不能解释为既有容量的无条件免费价值。

\subsection{代表节点缩放}

令某代表集团占行业服务份额$g$，拥有$F$个工厂，并定义单厂行业份额$q=g/F$。当前代表节点的AI需求比例$s^{\mathrm{AI}}_a$、基础负荷比例$s^{L}_a$和行业恢复倍数$M_a$为：

\begin{table}[htbp]
 \centering
 \caption{代表节点缩放与行业恢复。}
 \label{tab:representative-scaling}
 \begin{tabular}{lccc}
  \toprule
  架构 & $s^{\mathrm{AI}}_a$ & $s^L_a$ & $M_a$ \\
  \midrule
  IF & $q$ & $q$ & $F/g$ \\
  IG & $g$ & $q$ & $1/g$ \\
  II\_1host & $1$ & $q$ & $1$ \\
  \bottomrule
 \end{tabular}
\end{table}

三种架构均满足$s^{\mathrm{AI}}_aM_a=1$，因此恢复到行业层后提供相同有效服务。IG和II\_1host仍使用代表工厂基础负荷$qL_t$，而不是把全集团或全行业生产负荷搬到算力节点。\newwork 这一代表节点设计隔离了“AI算力集中程度”和“工业生产负荷空间分布”，但也意味着结果是代表桶筛查，不是具体地理节点的潮流结果。

\section{任务调度与CPU/GPU异构计算}

本节把逐时到达的有效服务转换为可执行的CPU/GPU工作量和服务器容量。模型要捕捉两个特点：一是柔性任务可以在截止期前移时，但不能消失；二是不同任务对CPU和GPU的适用性不同，因而同一服务量不必然对应同一GPU小时。我们用聚合批次和截止窗避免追踪每个请求，再用任务路由份额和成功吞吐换算硬件容量。

\subsection{聚合到达批次与截止窗}

\implemented 当前调度不排单个请求，也不使用显式积压状态。具有相同释放时刻$r_j$、截止窗$d_j$、柔性等级和任务ID的需求合并为到达批次$j$。循环代表周内允许执行小时集合为

\begin{equation}
 \mathcal T_j=\{(r_j+\tau)\bmod T:\tau=0,\ldots,d_j\},
 \label{eq:admissible-hours}
\end{equation}

执行量$x^a_{jt}$满足

\begin{equation}
 \sum_{t\in\mathcal T_j}x^a_{jt}=D_j,
 \qquad x^a_{jt}\geq0.
 \label{eq:job-conservation}
\end{equation}

当$d_j=0$时，任务只能在释放小时执行。周末释放的柔性任务可按循环稳态边界在周初允许时段完成。\boundary 该设计不表示真实跨周积压，也不是$O(T)$队列算法；变量规模随到达批次数和截止窗增长。

\literature 数据中心研究已证明延迟容忍计算可通过时间迁移改变功率形状。\adapted 本研究按制造业任务类别和截止窗形成聚合批次，并把服务守恒直接连接到异构服务器容量。

\subsection{任务到CPU/GPU计算需求}

这一转换不是按标称FLOPS比较CPU和GPU，而是询问在同样质量和SLA下，每种硬件完成一单位任务需要多少时间。路由份额表示任务类型的适用性，时间倍数表示选定硬件的成功吞吐差异。

令$h\in\Hset=\{\mathrm{gpu},\mathrm{cpu}\}$，$q_{kh}$为任务$k$路由到硬件$h$的份额，且$\sum_hq_{kh}=1$；$m_{kh}$为该硬件相对参考L20加速器小时的时间转换倍数；$a^{\mathrm{eff}}$为每单位有效服务所需参考加速器小时。时段$t$的硬件计算需求为

\begin{equation}
 H^a_{hnt}=a^{\mathrm{eff}}\sum_kq_{kh}m_{kh}d^a_{knt}.
 \label{eq:hardware-compute}
\end{equation}

若$G_h$为一个服务器组每小时可提供的服务容量，$N^{a,\mathrm{on}}_{hnt}$和$N^{a,\mathrm{inst}}_{nh}$分别为在线和安装服务器组，则

\begin{equation}
 H^a_{hnt}\leq G_hN^{a,\mathrm{on}}_{hnt},\qquad
 0\leq N^{a,\mathrm{on}}_{hnt}\leq N^{a,\mathrm{inst}}_{nh}.
 \label{eq:throughput}
\end{equation}

\boundary $q_{kh}$和$m_{kh}$必须由相同质量与SLA下的成功任务吞吐校准，不能仅由FLOPS、核心数或TDP推断。当前任务路由仍是证据约束的结构情景，而不是制造业总体观测比例。

\subsection{规划装机下限}

仅有平均吞吐约束会允许优化器过度依赖移时，从而给出缺乏工程裕量的服务器规模。因此，这一层希望捕捉“平常可以调度，但设备仍需能承受未错峰需求”的规划逻辑。\implemented 规划容量由未错峰到达时刻峰值预先计算，不由电价或调度结果反推。令$H^{a,\mathrm{arr}}_{hnt}$为按自然到达时刻计算的硬件需求，则

\begin{equation}
 N^{a,\mathrm{req}}_{nh}=\frac{\max_tH^{a,\mathrm{arr}}_{hnt}}{G_h},
 \label{eq:required-groups}
\end{equation}

\begin{equation}
 N^{a,\mathrm{floor}}_{nh}=
 \max\left\{(1+r_h^{\mathrm{plan}})N^{a,\mathrm{req}}_{nh},
 \left\lceil N^{a,\mathrm{req}}_{nh}\right\rceil+1\right\},
 \label{eq:planning-floor}
\end{equation}

\begin{equation}
 N^{a,\mathrm{inst}}_{nh}\geq N^{a,\mathrm{floor}}_{nh}.
 \label{eq:installed-floor}
\end{equation}

当前$r_h^{\mathrm{plan}}=10\%$，式~\eqref{eq:planning-floor}同时保留10\%规划裕量和N+1备用中较高者。优化变量可以是连续等效容量，但装机下限中的向上取整和加一仍保留小规模采购效应。正常运行不再逐时扣留额外10\%容量，已安装服务器可全部用于调度。

模型状态还要求GPU装机满足最小模型副本与自带存储：

\begin{equation}
 N^{a,\mathrm{inst}}_{n,\mathrm{gpu}}\geq N^{\mathrm{model,min}},\qquad
 S^{\mathrm{bundled}}N^{a,\mathrm{inst}}_{n,\mathrm{gpu}}
 \geq S^{\mathrm{model,req}}.
 \label{eq:model-state}
\end{equation}

这些下限用于表示备用、小规模采购和模型副本的最低物理要求。\boundary 它们尚未表示故障率、维修时间、多机通信、显存分片、容器启停或严格N+1可用性验证；10\%与加一规则仍是待敏感性检验的工程情景。

\section{三状态服务器功率与企业电力平衡}

服务器容量需要进一步转换为企业电表后的逐时负荷。本节希望捕捉“安装但冷备”、“在线但空闲”和“执行计算”三种状态的功率差异，以及AI负荷与基础生产负荷、屋顶光伏和电池的逐时叠加。这样，调度不仅改变计算完成时间，也会改变购电量和峰值。

\subsection{服务器和设施功率}

\implemented 当前模型不使用“安装台数乘综合利用率”的乘积式。令$P_h^{\mathrm{stby}}$、$P_h^{\mathrm{idle}}$和$P_h^{\max}$分别为冷备、在线空闲和满载整机功率，则设施功率为

\begin{equation}
\begin{aligned}
 P^{a,\mathrm{fac}}_{nt}=\sum_h\frac{\mathrm{PUE}_h}{10^3}\bigg[&
 P_h^{\mathrm{stby}}N^{a,\mathrm{inst}}_{nh}
 +(P_h^{\mathrm{idle}}-P_h^{\mathrm{stby}})N^{a,\mathrm{on}}_{hnt}\\
 &+\frac{P_h^{\max}-P_h^{\mathrm{idle}}}{G_h}H^a_{hnt}\bigg].
\end{aligned}
 \label{eq:facility-power}
\end{equation}

式~\eqref{eq:facility-power}依次表示：全部装机的冷备功率、在线设备从冷备升至空闲的增量，以及有效计算产生的动态功率。功率以kW输入并除以$10^3$换算为MW。配置中的硬件设施倍率已经包含冷却和供配电开销，不能在式外再次叠加冷却或内部网络系数。

\literature PUE的定义可依据ISO/IEC~30134-2；服务器空闲功率与负载相关动态功率也有数据中心测量研究支持。\adapted 将冷备、在线空闲和CPU/GPU计算量同时写入设施功率，是为当前架构比较构造的可审计线性近似。\pending 正式结果仍需目标CPU/GPU服务器的墙插功率--合格吞吐包络校准。

为单独观察算力池化，定义节点$n$的安装IT标称功率为

\begin{equation}
 K^{a,\mathrm{IT}}_n=\sum_hN^{a,\mathrm{inst}}_{nh}P_h^{\max}.
 \label{eq:installed-it-capacity}
\end{equation}

相对每家工厂独立配置的安装IT容量，架构$a$的池化收益为

\begin{equation}
 G^a_{\mathrm{pool}}
 =1-\frac{\sum_nK^{a,\mathrm{IT}}_n}
 {\sum_iK^{\mathrm{IT,stand}}_i}.
 \label{eq:pooling-gain}
\end{equation}

正值表示集中架构通过需求多样性和共享备用减少了安装IT标称容量。\boundary 该指标只表示设备层池化，不自动包含网络、治理、可用性、数据驻留或单点风险成本。

\subsection{光伏、储能与节点功率平衡}

节点功率平衡表示AI负荷不是独立用电，而是与工厂现有负荷和表后DER共同决定购电。光伏直接抵消同时负荷，电池在时间上转移净负荷；二者对AI的价值取决于AI任务的到达和截止窗，而不只取决于年用电量。

节点逐时电力平衡为

\begin{equation}
 p^a_{nt}+p^{a,\mathrm{PV,use}}_{nt}+p^{a,\mathrm{dis}}_{nt}
 =L^a_{nt}+P^{a,\mathrm{fac}}_{nt}+p^{a,\mathrm{ch}}_{nt},
 \label{eq:power-balance}
\end{equation}

其中$p^a_{nt}$为电网购电，$L^a_{nt}$为代表工厂基础负荷。储能采用固定时长，充放电效率各取往返效率平方根，并在完整168小时代表周边界循环：

\begin{equation}
 e^a_{n,t}=e^a_{n,t-1}
 +\eta_{\mathrm{ch}}p^{a,\mathrm{ch}}_{nt}\Delta t
 -\frac{p^{a,\mathrm{dis}}_{nt}\Delta t}{\eta_{\mathrm{dis}}},
 \qquad e^a_{n,0}=e^a_{n,T}.
 \label{eq:battery-balance}
\end{equation}

\boundary 无DER基准情景不配置光伏和电池。DER敏感性才允许既有屋顶光伏和电池投资；不能把敏感性配置描述为所有核心情景的共同条件。

\section{无AI反事实与新增接入容量}

本节要从工厂本来就需要的用电中分离出AI带来的额外接入需求。直接使用AI情景的总峰值会把基础生产负荷也归因给AI；直接假定大量未观测裕量又会低估AI的容量影响。因此，每个AI情景都与价格、基础负荷和DER规则相同的无AI反事实成对比较，并只报告净峰值超出基线的正部分。

\subsection{匹配的无AI基线}

\implemented 每个代表节点先在与AI情景相同的价格、基础负荷和DER规则下求解无AI反事实。既有接入边界取其优化后净购电峰值：

\begin{equation}
 K^{0}_n=\max_t p^{\mathrm{noAI}}_{nt}.
 \label{eq:baseline-capacity}
\end{equation}

AI情景的新增接入容量为

\begin{equation}
 \Delta K^a_n=\left[\max_t p^a_{nt}-K^0_n\right]^+,
 \qquad[z]^+=\max(0,z).
 \label{eq:grid-expansion}
\end{equation}

该式不给未观测合同容量、配变、馈线或变电站余量任何信用。由于无AI与AI情景采用相同DER规则，基线DER造成的净峰值降低也不会被误认为AI可免费使用的余量。

\newwork “优化无AI净峰值作为匹配接入边界”是本研究为防止免费余量信用而采用的反事实设计。它是代表节点容量筛查，不等于真实企业已签合同容量或实际配网可用容量。

\subsection{容量工程指标}

行业层至少报告总新增接入容量和最大代表节点接入容量：

\begin{equation}
 \Delta K^a_{\Sigma}=\sum_n\Delta K^a_n,
 \qquad
 \Delta K^a_{\max}=\max_n\Delta K^a_n.
 \label{eq:grid-indicators}
\end{equation}

前者描述代表接入容量总量，后者描述最大单点连接规模。\boundary 二者均不能在没有资产、阈值和项目时序数据时直接称为已避免的具体工程。

\subsection{相对参考架构避免的接入容量}

以共同参考架构$a_0$为基准，架构$a$避免的新增接入容量及其比例定义为

\begin{align}
 \Delta K^{\mathrm{avoid}}_{a\mid a_0}
 &=\Delta K^{a_0}_{\Sigma}-\Delta K^a_{\Sigma},
 \label{eq:avoided-capacity}\\
 R^{\mathrm{avoid}}_{a\mid a_0}
 &=\frac{\Delta K^{a_0}_{\Sigma}-\Delta K^a_{\Sigma}}
 {\Delta K^{a_0}_{\Sigma}},
 \qquad \Delta K^{a_0}_{\Sigma}>0.
 \label{eq:avoided-capacity-fraction}
\end{align}

\implemented 式~\eqref{eq:avoided-capacity}的单位是MW，式~\eqref{eq:avoided-capacity-fraction}为无量纲比例。它们只表示不同部署架构下模型接入容量需求的差异。当前研究不估计电网升级的货币成本，不使用单位容量造价、项目实现系数或延期价值，也不把避免的MW称为“避免的容量投资”。

\boundary 新增接入容量是连续物理筛查指标，不等于已批复的配变、馈线或变电站工程。在没有离散设备档位、现有裕量、线路长度、工程时序和地区造价时，不由容量差推断具体投资节省。

\section{企业购电成本与当前年度目标}

这一节把物理运行结果转换为企业决策口径的年度成本。希望捕捉的是：任务移时既可能改变何时买电，也可能改变月度最大需量；分散部署可能重复购买服务器和运维多个节点，集中部署则可能获得池化收益。这些成本在共同年份和企业付款边界下加总。新增接入容量仍是独立物理结果，不在本节折算为电网升级成本。

\subsection{到户电量费与最大需量费}

企业电费同时取决于用了多少电和形成多大计费需量。因此，电量价按逐时购电累积，最大需量另行按峰值计费，而不把需量费平均摊入每度电。\implemented 当前小时价格是广东日前代表周价格加体积附加：

\begin{equation}
 \pi_t^{\mathrm{retail}}
 =\pi_t^{\mathrm{DA}}+10^3
 (c^{\mathrm{loss}}+c^{\mathrm{TD}}+c^{\mathrm{sys}}+c^{\mathrm{gov}}).
 \label{eq:retail-price}
\end{equation}

代表周权重为$w^{\mathrm{ann}}=365/7$。电量费与最大需量费为

\begin{align}
 C^a_{\mathrm{energy}}
 &=w^{\mathrm{ann}}\sum_{n,t}\pi_t^{\mathrm{retail}}p^a_{nt}\Delta t,
 \label{eq:energy-cost}\\
 C^a_{\mathrm{demand}}
 &=12\times10^3c^{\mathrm{dem}}\sum_n\max_tp^a_{nt}.
 \label{eq:demand-cost}
\end{align}

\boundary 这些是企业付款代理，不是社会边际发电成本，也不代表全国制造业实际合同组合。相对无AI的增量必须由两个完整反事实成本之差计算。

\literature SIELS研究采用外生逐时边际电能成本和发电容量成本评价工业负荷灵活性，并明确排除输配电延期价值。\adapted 若未来获得适用的社会边际成本，可以把架构间有符号逐时负荷差作为独立后处理。该批发侧价值不能与式~\eqref{eq:energy-cost}的企业电量费相加，否则会重复计算同一电量价值。

\subsection{服务器与储能年化成本}

设备采购是一次性投资，而企业电费和运维是年度支出。为了在同一目标函数中比较，本节将服务器和电池投资按各自寿命年化，再加入固定运维。

贴现率为$r$、寿命为$Y$的企业资产使用资本回收因子

\begin{equation}
 \operatorname{CRF}(r,Y)=\frac{r(1+r)^Y}{(1+r)^Y-1}.
 \label{eq:crf}
\end{equation}

服务器类型$h$的采购价为$I_h$，设施资本附加比例为$f_h^{\mathrm{fac}}$，寿命为$Y_h$，年度维护比例为$f_h^{\mathrm{OM}}$，则每组年化成本为

\begin{equation}
 c_h^{\mathrm{server}}
 =I_h(1+f_h^{\mathrm{fac}})\operatorname{CRF}(r,Y_h)
 +I_hf_h^{\mathrm{OM}}.
 \label{eq:server-cost}
\end{equation}

对固定时长$d_B$的电池，能量投资$c_E$、逆变器功率投资$c_P$、相应寿命$Y_E,Y_P$和逆变器固定运维$f_P^{\mathrm{FOM}}$给出

\begin{equation}
 c^B=10^3d_Bc_E\operatorname{CRF}(r,Y_E)
 +10^3c_P[\operatorname{CRF}(r,Y_P)+f_P^{\mathrm{FOM}}].
 \label{eq:battery-cost}
\end{equation}

\subsection{八项增量成本目标}

总目标将设备、运行和部署点生命周期账户纳入同一年度口径。这样可以比较“多节点重复成本”和“集中池化收益”，而不只比较电费或服务器采购中的某一项。\implemented 企业侧联合优化层的年成本为

\begin{equation}
\begin{aligned}
 J_a^{\mathrm{net}}={}&
 \sum_{n,h}c_h^{\mathrm{server}}N^a_{nh}
 +c^{\mathrm{PV}}K^a_{\mathrm{PV}}
 +c^BK^a_B\\
 &+C^a_{\mathrm{energy}}+C^a_{\mathrm{demand}}.
\end{aligned}
 \label{eq:network-objective}
\end{equation}

模型初始化、模型增量存储和按部署点运维作为年度生命周期项加入：

\begin{equation}
 J_a=J_a^{\mathrm{net}}+C_a^{\mathrm{init}}
 +C_a^{\mathrm{storage}}+C_a^{\mathrm{operations}}.
 \label{eq:total-cost}
\end{equation}

行业$j$的架构等效增量成本由匹配的无AI代表节点恢复：

\begin{equation}
 \Delta J^a_j=M^a_j\sum_{c\in\mathcal C}
 (C^a_{jc}-C^{\mathrm{noAI}}_{jc}),
 \label{eq:incremental-cost}
\end{equation}

其中成本集合$\mathcal C$恰好包括八项：服务器、PV、电池、电量、最大需量、模型初始化、模型存储和模型运维。新增接入容量由式~\eqref{eq:grid-expansion}--\eqref{eq:avoided-capacity-fraction}单独报告，不属于$\mathcal C$，也不进入$J_a^{\mathrm{net}}$。

\boundary 当前核心目标不包含电网升级成本、网络传输、碳、水、土地、材料、噪声、未服务任务或风险惩罚。它们可以作为独立物理指标、后处理或扩展目标，但不能写成已经进入当前$J_a$。

\section{企业计算采购方式：纯自建、纯云端与混合部署}

本节在部署尺度$a\in\{\mathrm{IF},\mathrm{IG},\mathrm{II}_{1\mathrm{host}}\}$之外，再区分企业如何取得计算服务。部署尺度回答“自建设备放在多大的代表节点”，采购方式回答“任务由自建设备还是云端预留容量完成”。代码把两层连接在同一个逐时优化中，但为避免把成本边界混在一起，下面分别写出纯自建、纯云端和混合部署三个模型。纯自建是现有核心模型；纯云端和混合部署是已实现的单行业结构敏感性，不等同于另一个改变生产负荷位置的物理架构。

\subsection{模型一：纯自建}

纯自建要求全部有效服务由企业自有CPU/GPU设备完成。对柔性批次$j$，本地执行量$x^a_{jt}$仍满足式~\eqref{eq:job-conservation}；刚性服务$R^a_{knt}$按到达时刻留在本地。其企业目标就是式~\eqref{eq:total-cost}：

\begin{equation}
 J^{\mathrm{own}}_{a,d}=J_{a,d},
 \label{eq:pure-owned-model}
\end{equation}

并保留式~\eqref{eq:planning-floor}--\eqref{eq:model-state}的未错峰装机下限、最低模型副本和自带存储约束。纯自建模型没有云订阅成本。

\subsection{模型二：纯云端预留容量}

当前代码中的纯云端边界不是Token API或SaaS，而是企业购买年度CPU/GPU预留容量。令$y^a_{jt}$为柔性批次在云端的执行量，$z^a_{knt}$为刚性任务在到达时刻转交云端的服务量；纯云端固定$x^a_{jt}=0$和本地刚性执行量$R^a_{knt}-z^a_{knt}=0$。云端硬件$h$的逐时计算需求为

\begin{equation}
 H^{a,\mathrm{cloud}}_{hnt}
 =a^{\mathrm{eff}}\sum_kq_{kh}m_{kh}d^{a,\mathrm{cloud}}_{knt},
 \label{eq:cloud-compute}
\end{equation}

并由年度预留组数$B^a_{nh}$满足

\begin{equation}
 H^{a,\mathrm{cloud}}_{hnt}
 \leq G_h^{\mathrm{cloud}}B^a_{nh},
 \qquad B^a_{nh}\geq0.
 \label{eq:cloud-reserved-capacity}
\end{equation}

若$c_h^{\mathrm{sub}}$为每组云容量的年度订阅价，则当前代码对应的纯云端目标为

\begin{equation}
 J^{\mathrm{cloud}}_{a,d}
 =J_{0,d}+\sum_{n,h}c_h^{\mathrm{sub}}B^a_{nh}
 +C_a^{\mathrm{init}}+C_a^{\mathrm{operations}}.
 \label{eq:pure-cloud-model}
\end{equation}

在全云解中，本地AI服务器、AI设施用电和模型增量存储为零；工厂基础负荷及其DER仍存在，但与相同能源边界的无AI反事实相减。\boundary 代码仍对全云解计入每部署点的模型初始化和运维成本；云端对象存储、流量、迁移集成、支持和税费没有进入这个联合优化模型。

\subsection{模型三：自建与云端联合优化}

混合模型允许每项柔性任务在截止窗内分别由本地或云端执行，并允许每项刚性任务在原到达时刻分配到两侧。服务守恒分别写为

\begin{align}
 \sum_{t\in\mathcal T_j}\left(x^a_{jt}+y^a_{jt}\right)&=D_j,
 \label{eq:hybrid-flexible-conservation}\\
 d^{a,\mathrm{local}}_{knt}+d^{a,\mathrm{cloud}}_{knt}
 &=R^a_{knt},
 \qquad 0\leq d^{a,\mathrm{cloud}}_{knt}\leq R^a_{knt}.
 \label{eq:hybrid-rigid-conservation}
\end{align}

云端份额还受配置上限$\bar s^{\mathrm{cloud}}$约束：

\begin{equation}
 \sum_{n,k,t}d^{a,\mathrm{cloud}}_{knt}
 \leq\bar s^{\mathrm{cloud}}
 \sum_{n,k,t}\left(d^{a,\mathrm{local}}_{knt}
 +d^{a,\mathrm{cloud}}_{knt}\right).
 \label{eq:maximum-cloud-share}
\end{equation}

本地任务继续进入CPU/GPU装机、在线状态、三状态功率、购电、DER和最大需量模型；云端任务进入式~\eqref{eq:cloud-compute}和\eqref{eq:cloud-reserved-capacity}，不进入工厂电力平衡。当前联合目标为

\begin{equation}
\begin{aligned}
 J^{\mathrm{hyb}}_{a,d}={}&J^{\mathrm{local}}_{a,d}
 +\sum_{n,h}c_h^{\mathrm{sub}}B^a_{nh}\\
 &+C_a^{\mathrm{init}}+C_a^{\mathrm{storage}}
 +C_a^{\mathrm{operations}}+\phi^{\mathrm{grid}}K_a^{\mathrm{exp}}.
\end{aligned}
 \label{eq:hybrid-model}
\end{equation}

其中$J^{\mathrm{local}}_{a,d}$包含本地服务器、PV、电池、电量费和最大需量费；$\phi^{\mathrm{grid}}K_a^{\mathrm{exp}}$只在“高增容惩罚”结构压力测试中启用，核心值为零。它是迫使模型在本地增容、储能和云订阅之间显露替代关系的数值稀缺惩罚，不是电网升级成本，也不得作为企业成本参数解释。

式~\eqref{eq:hybrid-model}把前两个模型的成本机制串联起来：当$\bar s^{\mathrm{cloud}}=1$且最优本地执行为零时，它退化为当前预留容量口径的纯云端；本地和云端执行均为正的内生解才称为混合部署。需要注意，$\bar s^{\mathrm{cloud}}=0$只会使服务分配变为全本地，并不会严格恢复模型一，因为代码在混合开关打开时仍取消了部分本地固定下限。因此纯自建必须作为独立模型求解，不能用混合模型的零云份额替代。企业可在三个独立求解结果之间比较

\begin{equation}
 \Delta C^{\mathrm{ent}}_{a,d,b}
 =M_a\left(J^{*}_{a,d,b}-J^{*}_{0,d}\right),
 \qquad b\in\{\mathrm{own},\mathrm{cloud},\mathrm{hyb}\},
 \label{eq:procurement-incremental-cost}
\end{equation}

并选择

\begin{equation}
 b_{a,d}^{*}=\arg\min_{b\in
 \{\mathrm{own},\mathrm{cloud},\mathrm{hyb}\}}
 \Delta C^{\mathrm{ent}}_{a,d,b}.
 \label{eq:procurement-choice}
\end{equation}

\implemented 当前代码已经联合决定柔性和刚性任务的本地/云端分配、本地CPU/GPU装机与运行、云端CPU/GPU预留组数、电池和购电，并验证总服务量守恒。\newwork 需要作者重点判断两项实现边界。第一，混合开关打开后，当前实现取消外部“10\%与N+1取较高值”的装机下限、最低在线GPU以及本地模型副本/存储下限；本地装机仍受按实际调度峰值计算的10\%规划裕量约束。这使纯云解可行，但也可能低估“仍保留少量本地任务”时的最小现场固定成本。第二，代码目前只限制云端总服务份额，没有按任务应用式~\eqref{eq:task-admissibility}的数据驻留矩阵；因此隐私或硬实时任务在数学上也可被分给云端。正式采用混合模型前，应决定是否加入任务级云端可接受性，以及条件触发的本地副本、N+1和运维固定成本。

\boundary 上述联合模型使用同一组$q_{kh}$和$m_{kh}$将本地与云端服务换算为CPU/GPU需求，并按公开价格代理购买预留容量。它没有表示云厂商内部电力、服务器或接入容量，不能把这些供应商资源成本再加到企业订阅费上；也没有验证云端与本地硬件在共同质量和SLA下的真实等效吞吐。

\section{企业部署决策：无DER与有DER两种边界}

前述三个采购模型给出每种部署尺度下的企业成本，本节进一步比较能源边界。已有或可投资的光伏和储能只直接作用于本地执行份额，因此既可能改变部署尺度，也可能改变本地/云端分配；无DER与有DER不能共用一个反事实。当前核心架构比较仍使用纯自建成本，混合敏感性则用式~\eqref{eq:procurement-incremental-cost}在同一能源边界内比较采购组合。

\subsection{共同决策原则}

令$d\in\{0,1\}$表示企业能源边界：$d=0$为无分布式发电和储能，$d=1$为允许分布式光伏和储能。对每个架构$a$和能源边界$d$，企业先求解

\begin{equation}
 J_{a,d}^{*}=\min_{\Omega_{a,d}}J_{a,d},
 \label{eq:enterprise-solve}
\end{equation}

其中$\Omega_{a,d}$包含相同有效服务、SLA、CPU/GPU吞吐、装机下限、任务调度和接入容量约束，并按$d$改变DER可行域。相同能源边界下的无AI反事实成本记为$J_{0,d}^{*}$。行业等效企业增量成本为

\begin{equation}
 \Delta C^{\mathrm{ent}}_{a,d}
 =M_a\left(J_{a,d}^{*}-J_{0,d}^{*}\right).
 \label{eq:enterprise-incremental-cost}
\end{equation}

\implemented 企业在满足共同服务与数据边界的可行架构集合$\Aset^{\mathrm{feas}}$中选择增量成本最低者：

\begin{equation}
 a_d^{*}=\arg\min_{a\in\Aset^{\mathrm{feas}}}
 \Delta C^{\mathrm{ent}}_{a,d}.
 \label{eq:enterprise-choice}
\end{equation}

相对参考架构$a_0$的企业成本优势定义为

\begin{equation}
 V^{\mathrm{ent}}_{a\mid a_0,d}
 =\Delta C^{\mathrm{ent}}_{a_0,d}
 -\Delta C^{\mathrm{ent}}_{a,d}.
 \label{eq:enterprise-value}
\end{equation}

当$V^{\mathrm{ent}}_{a\mid a_0,d}>0$时，架构$a$在能源边界$d$下比参考架构具有较低企业增量成本。隐私、质量、数据驻留和硬时延首先作为可行性条件，不在当前核心决策式中货币化。

\subsection{情景一：无分布式发电和储能}

无DER情景是识别计算部署本身的基准。它将光伏和电池固定为零，使架构差异只能通过服务器装机、任务调度、基础负荷的非同时性和企业电价体现。

无DER可行域固定

\begin{equation}
 K^{a}_{\mathrm{PV}}=K^a_B=0,\qquad
 p^{a,\mathrm{PV,use}}_{nt}=p^{a,\mathrm{ch}}_{nt}
 =p^{a,\mathrm{dis}}_{nt}=0.
 \label{eq:no-der-domain}
\end{equation}

因此$p^a_{nt}=L^a_{nt}+P^{a,\mathrm{fac}}_{nt}$，无AI基线为$p^{\mathrm{noAI}}_{nt}=L^a_{nt}$。该边界下企业目标为

\begin{equation}
\begin{aligned}
 J_{a,0}={}&C_a^{\mathrm{server}}
 +C_{a,0}^{\mathrm{energy}}+C_{a,0}^{\mathrm{demand}}\\
 &+C_a^{\mathrm{init}}+C_a^{\mathrm{storage}}
 +C_a^{\mathrm{operations}}.
\end{aligned}
 \label{eq:enterprise-no-der}
\end{equation}

式~\eqref{eq:enterprise-no-der}隔离计算部署、任务调度和生产负荷非同时性对企业电费的作用，不给予光伏或储能任何贡献。新增接入容量仍由式~\eqref{eq:grid-expansion}计算，但不进入企业目标函数。

\subsection{情景二：允许分布式光伏和储能}

有DER情景要检验工厂现场能源资源是否改变AI部署的边际价值。光伏可与白天AI和生产负荷同时消纳，电池可在时段之间转移净负荷；但它们也可能在无AI时就已对基础负荷有价值。因此，本情景必须与拥有相同DER选择的无AI基线比较。

有DER可行域加入式~\eqref{eq:power-balance}和\eqref{eq:battery-balance}，并施加

\begin{equation}
 0\leq K^a_{\mathrm{PV}}\leq\overline K^a_{\mathrm{PV}},
 \qquad K^a_B\geq0,
 \label{eq:der-domain}
\end{equation}

以及逐时光伏出力和电池功率/能量上限。有DER企业目标为

\begin{equation}
\begin{aligned}
 J_{a,1}={}&C_a^{\mathrm{server}}
 +C_{a,1}^{\mathrm{energy}}+C_{a,1}^{\mathrm{demand}}\\
 &+C_{a,1}^{\mathrm{PV}}+C_{a,1}^{B}
 +C_a^{\mathrm{init}}+C_a^{\mathrm{storage}}
 +C_a^{\mathrm{operations}}.
\end{aligned}
 \label{eq:enterprise-with-der}
\end{equation}

当前DER敏感性把屋顶光伏视为已有资源：$K^a_{\mathrm{PV}}$固定在情景给定上限且$C_{a,1}^{\mathrm{PV}}=0$，同时允许电池投资。若未来研究企业新建光伏，则$K^a_{\mathrm{PV}}$应成为决策变量，并按年化投资计入$C_{a,1}^{\mathrm{PV}}$。

\boundary 有DER时，无AI反事实$J_{0,1}^{*}$必须拥有与AI情景相同的屋顶、储能投资选择、价格和运行规则。不得用无DER基线$J_{0,0}^{*}$扣减有DER的AI成本，否则会把基础负荷本来可以获得的DER收益错误归因于AI。

\subsection{DER对AI增量成本的净贡献}

DER对架构$a$的净价值用两个能源边界下的企业增量成本之差表示：

\begin{equation}
 V^{\mathrm{DER}}_a
 =\Delta C^{\mathrm{ent}}_{a,0}
 -\Delta C^{\mathrm{ent}}_{a,1}.
 \label{eq:der-difference-in-differences}
\end{equation}

式~\eqref{eq:der-difference-in-differences}等价于

\begin{equation}
 V^{\mathrm{DER}}_a
 =M_a\left[(J_{a,0}^{*}-J_{0,0}^{*})
 -(J_{a,1}^{*}-J_{0,1}^{*})\right].
 \label{eq:der-did-expanded}
\end{equation}

正值表示DER降低了AI带来的增量企业成本；负值表示在当前投资、价格和运行边界下，新增DER成本超过其电量费和最大需量费收益。DER对新增接入容量的影响另作物理指标报告，不进入该货币价值。\newwork 这一双重差分是企业决策解释式，不是独立的物理优化变量；它从四个匹配反事实的求解结果重构。

DER是否改变架构排序，可比较

\begin{equation}
 \left(\Delta C^{\mathrm{ent}}_{a,0}-\Delta C^{\mathrm{ent}}_{b,0}\right)
 \left(\Delta C^{\mathrm{ent}}_{a,1}-\Delta C^{\mathrm{ent}}_{b,1}\right)<0.
 \label{eq:der-ranking-switch}
\end{equation}

若式~\eqref{eq:der-ranking-switch}成立，说明加入光伏和储能后架构$a$与$b$的企业成本排序发生反转；这比只报告某一情景下的总成本更直接回答企业部署决策问题。

\section{敏感性分析：结论驱动因素与架构排序}
\label{sec:sensitivity-analysis}

本节集中回答两个不同问题：第一，哪些输入决定AI用电、服务器数量、接入容量和企业成本的\emph{绝对总量}；第二，哪些输入会改变工厂侧、集团池、行业节点或云服务之间的\emph{相对成本与排序}。二者不能用同一判据判断。某参数使全国AI成本增加一倍，并不自动表示最便宜的架构发生变化；反之，一个对总量影响不大的固定成本或价格差，也可能使两个成本接近的架构发生排序反转。

\pending 当前全国31行业联合CPU/GPU结果尚未完整重跑，C38单行业第一轮筛查也主要检验物理输出摆幅，不能据此声称某一架构在全国范围内更便宜或成本排序已经稳健。因此，本节先冻结分析问题、参数分组、比较判据和结果表位置；所有数值、临界值和稳健性等级均待同版本测试完成后填写。

\subsection{总需求何时只改变总量，何时也会改变排序}

令制造业AI总需求尺度为$\lambda$，任务组合为$\bm\pi=(\pi_1,\ldots,\pi_K)$，其中$\sum_k\pi_k=1$。需求可以写为

\begin{equation}
 D_k(\lambda,\bm\pi)=\lambda\pi_k\bar D,
 \label{eq:sensitivity-demand-decomposition}
\end{equation}

其中$\bar D$是共同参考服务量。$\lambda$回答“总共需要多少AI服务”，$\bm\pi$回答“这些服务由哪些任务组成”。若所有架构的服务需求同比例变化，任务组合、SLA、CPU/GPU路由和价格均不变，且设备、云容量与接入容量均可连续缩放，则企业增量成本可近似表示为

\begin{equation}
 \Delta C_a(\lambda,\bm\pi)
 \approx F_a+\lambda c_a(\bm\pi),
 \label{eq:sensitivity-cost-scale}
\end{equation}

其中$F_a$为模型副本、最低运维、N+1备用、最小机房或最低订阅等固定成本，$c_a$为单位有效服务的边际成本。两个架构$a$和$b$的成本差为

\begin{equation}
 \Delta C_a-\Delta C_b
 =(F_a-F_b)+\lambda\left[c_a(\bm\pi)-c_b(\bm\pi)\right].
 \label{eq:sensitivity-ranking-margin}
\end{equation}

因此，用户提出的直觉在一个受限条件下成立：当$F_a-F_b$可忽略、所有容量近似连续、没有价格分档或资源约束，而且$\bm\pi$不变时，共同改变$\lambda$主要按比例改变服务器、用电和成本总量，单位服务成本及架构排序应近似不变。此时应同时报告总成本和单位有效服务成本，以把规模效应与架构效应分开。

但本研究不能预先假定该比例不变性。工厂侧部署存在按站点重复的模型状态、运维和N+1备用，服务器采购可能整数化；集团池和行业节点存在池化门槛；云服务可能按预留容量、最低承诺或阶梯折扣计价；接入容量、DER和最大需量也受峰值与约束是否绑定影响。只要存在这些固定项、离散项或容量门槛，改变$\lambda$就可能改变单位成本并使式~\eqref{eq:sensitivity-ranking-margin}过零。若两架构的单位成本斜率不同，连续近似下的需求打平点为

\begin{equation}
 \lambda^*_{a,b}
 =-\frac{F_a-F_b}{c_a(\bm\pi)-c_b(\bm\pi)},
 \label{eq:sensitivity-demand-break-even}
\end{equation}

但只有在打平点附近成本近似线性且分母非零时才报告该值；含整数定容或分段价格时，应直接通过情景求解寻找排序变化区间。

\subsection{决定“哪个更便宜”的参数层级}

表~\ref{tab:sensitivity-driver-map}按参数是否主要改变总量、单位成本或可行域，将现有模型中的敏感因素归为四组。该分类是测试假设，不是结果；最终仍需用成本差、比值和临界值验证。

\begin{longtable}{p{2.6cm}p{4.7cm}p{3.6cm}p{3.0cm}}
 \caption{敏感因素、作用机制与需要保护的结论。}\label{tab:sensitivity-driver-map}\\
 \toprule
 因素组 & 代表参数或结构情景 & 主要影响路径 & 对成本排序的预期作用 \\
 \midrule
 \endfirsthead
 \toprule
 因素组 & 代表参数或结构情景 & 主要影响路径 & 对成本排序的预期作用 \\
 \midrule
 \endhead
 共同规模 & AI总服务需求$\lambda$、所有任务共同的调用强度、共同服务--计算效率 & 改变服务器、用电、容量和成本的绝对规模 & 在线性连续区间内应主要改变总量；必须检验固定成本、整数定容和容量门槛导致的反转 \\
 任务构成与约束 & 各任务份额$\bm\pi$、刚性/柔性比例、允许窗口、CPU/GPU路由、数据驻留、质量、时延和可靠性要求 & 改变峰值、池化机会、硬件组合和可选执行地点 & 可直接改变相对成本；不能用总需求统一缩放替代 \\
 本地与云价格 & GPU/CPU服务器价格、寿命、维护和设施附加；电量与最大需量电价；云API或预留容量价格、容量等效系数、合同折扣、存储与网络费 & 改变本地固定/边际成本与云端付款斜率 & 是“自建还是云服务更便宜”的一级直接驱动，应报告价格打平点 \\
 部署结构与能源边界 & 工厂/集团节点数、集团份额、N+1、装机裕量、整数采购、最低本地副本与运维、DER、接入上限和本地--云端可行性 & 改变重复固定成本、统计复用、峰值和可行解集合 & 可能产生非连续反转，应作为结构敏感性而非普通弹性报告 \\
 环境后处理 & WUE、电网排放因子、工厂建筑复用率、云端新建比例、建筑和材料强度 & 改变水、碳、土地和材料结果，不进入当前八项企业成本 & 不应改变“企业付款谁更低”，但会改变社会与环境权衡结论 \\
 \bottomrule
\end{longtable}

任务占比$\bm\pi$尤其重要。办公与Agent任务更可能使用Token API或CPU/云服务，预测维护、排程和仿真对CPU/GPU路由与批处理窗口更敏感，数据驻留和硬实时任务则可能排除部分云端方案。因而即使总服务量$\lambda\bar D$完全相同，行业或企业的任务组合不同，也可能对应不同的最便宜架构。正式分析应将“行业AI需求规模”与“行业任务构成”分开，不以某行业总成本高就推断其单位服务成本或某种架构更差。

\subsection{需要完成的测试与报告判据}

敏感性分析按以下顺序完成。每项测试均保持相同有效服务、质量和SLA，不允许通过减少服务量制造成本优势。

\begin{enumerate}
  \item \textbf{需求尺度测试：}在低、基准和高需求情景下保持$\bm\pi$不变，重新求解各架构。报告绝对成本、单位有效服务成本、成本比、排序和新增接入容量；若排序不变，再检验成本是否近似同比例变化，而不能仅凭方向相同宣称尺度无关。
  \item \textbf{任务构成测试：}在保持总服务量不变的条件下改变办公、Agent、VLM、维护、排程和仿真的份额，并同步应用对应的CPU/GPU路由、柔性和数据边界。优先比较代表性任务组合，不将六类任务的高值彼此独立相乘。
  \item \textbf{价格与采购测试：}分别改变本地GPU/CPU服务器采购价、寿命、维护、设施附加、电价和最大需量费，以及云端API/预留容量价格、企业交付倍数和容量等效系数。对每个直接价格驱动求使$\Delta C_a=\Delta C_b$的临界值，并与证据支持范围比较。
  \item \textbf{部署结构测试：}检验连续与整数服务器组、0与N+1备用、0--25\%装机裕量、条件触发的本地模型副本和运维、集团规模/工厂数、无DER与有DER，以及任务级云端可接受性。该组报告离散情景和排序反转，不报告伪精确弹性。
  \item \textbf{联合不利情景：}对最可能缩小成本优势的3--5个因素运行预先登记的联合角点。只有单因素和联合不利情景均不翻转，才可把结论称为范围稳健。
  \item \textbf{行业异质性检查：}全国核心结果完成后，同时比较行业年总成本与单位有效服务成本。前者说明需求规模，后者用于判断任务结构、节点数量、路由和价格如何影响相对经济性；描述性相关不得写成因果归因。
\end{enumerate}

对任意两个架构$a$和$b$，主要排序指标为

\begin{equation}
 M^{\mathrm{cost}}_{a,b}(\theta)
 =\frac{\Delta C_b(\theta)-\Delta C_a(\theta)}{\Delta C_b(\theta)},
 \label{eq:sensitivity-cost-margin}
\end{equation}

其中$M^{\mathrm{cost}}_{a,b}>0$表示$a$比$b$便宜。敏感性结果不能只报告某一输出相对基准变化了多少，而应优先报告最小成本余量、是否反转以及临界值。稳健性分为：证据范围和联合不利角点内均不反转的“范围稳健”；存在透明成立阈值的“条件稳健”；证据范围覆盖翻转点的“不稳健”。

\small
\begin{longtable}{p{1.5cm}p{2.5cm}p{2.1cm}p{2.0cm}p{1.8cm}p{1.5cm}p{1.7cm}}
 \caption{架构成本排序敏感性结果预留表。}\label{tab:sensitivity-ranking-placeholder}\\
 \toprule
 因素ID & 测试范围/情景 & 比较架构 & 最小成本余量 & 排序是否反转 & 临界值 & 当前状态 \\
 \midrule
 \endfirsthead
 \toprule
 因素ID & 测试范围/情景 & 比较架构 & 最小成本余量 & 排序是否反转 & 临界值 & 当前状态 \\
 \midrule
 \endhead
 PHY01 & 需求低/基准/高 & IF、IG、II、云 & 待计算 & 待计算 & 待计算 & \pending \\
 PHY03、HW01--02 & 柔性、路由和CPU相对服务时间 & IF、IG、II、云 & 待计算 & 待计算 & 待计算 & \pending \\
 COST组、HW03 & 本地与云端价格、寿命和容量等效 & 自建与云 & 待计算 & 待计算 & 待计算 & \pending \\
 PHY07--09 & 装机裕量、N+1和整数定容 & IF、IG、II & 待计算 & 待计算 & 不一定连续 & \pending \\
 PHY12 & 无DER/有DER & IF、IG、II、混合 & 待计算 & 待计算 & 结构情景 & \pending \\
 JOINT & 联合不利角点 & 由单因素筛选 & 待计算 & 待计算 & 待计算 & \pending \\
 \bottomrule
\end{longtable}
\normalsize

\subsection{当前可陈述的结论与尚不能陈述的结论}

现阶段可以确定的是分析逻辑，而不是最终架构胜负：（1）AI总需求规模必然影响服务器、用电、接入容量和成本总量；（2）只有在近似线性、连续且任务组合固定的区间内，总需求共同缩放才预期不改变相对单位成本和排序；（3）任务组合、CPU/GPU路由、可移时性、数据边界以及本地与云价格直接决定“哪个更便宜”；（4）固定成本、站点数量、N+1和整数采购使小规模企业最可能偏离比例关系；（5）环境足迹参数不应被混入企业付款排序。

\pending 现阶段不能陈述“AI总需求不会影响哪个架构更便宜”，也不能陈述“某架构普遍最便宜”。可以检验的、更严格的命题是：在预先登记的需求区间内，保持任务组合和其他条件不变时，架构排序是否保持一致；若不一致，则报告需求打平区间及导致反转的固定成本或容量门槛。全国联合结果、完整价格敏感性和联合不利角点完成后，再用表~\ref{tab:sensitivity-ranking-placeholder}替换这些条件性表述。

\section{建议报告的结果结构}

结果报告应沿着“服务是否等价---如何转换成设备和功率---如何影响容量和成本---结论对假设是否稳健”的链条展开。这个顺序先证明架构完成了同样服务，再解释差异来自哪一层机制，避免直接从总成本或容量排序跳到因果结论。

全国联合异构重跑完成后，每个行业和架构应按相同顺序报告：

\begin{enumerate}
  \item 有效服务总量、刚性/柔性任务完成情况和SLA；
  \item CPU/GPU计算需求、安装服务器组和逐时在线服务器组；
  \item 设施峰值、年设施用电及三状态功率分解；
  \item 代表节点新增接入容量、行业等效容量总和和最大单点容量；
  \item 相对参考架构避免的新增接入容量（MW）和避免比例；
  \item 企业电量费、最大需量费和八项增量成本分解；
  \item 混合敏感性中的本地/云端任务份额、云端GPU/CPU预留组数、云订阅费和压力测试惩罚项；
  \item 无DER和有DER两种边界下的企业增量成本、架构选择与DER双重差分价值；
  \item 任务调度、CPU路由、装机下限、功率参数和DER配置敏感性。
\end{enumerate}

\pending 当前只确认C36--IF单案例在服务守恒、CPU/GPU功率、服务器成本和总成本复算上通过。全国31行业$\times\{\mathrm{IF},\mathrm{IG},\mathrm{II}_{1\mathrm{host}}\}$尚未完成同口径重跑，因此本版不提供架构排序、全国节省比例或中美比较。

\section{方法来源、原创性与需作者判断的内容}

本节不再引入新的数学模型，而是把前述方法按“直接采用的文献机制、为本问题改写的机制、已进入实现的定义、以及本研究新增的选择”重新整理。目的是让审稿人能够区分通用方法与本文的实质性建模判断，也提醒作者哪些内容尚需外部验证或敏感性分析。

\begin{longtable}{p{4.0cm}p{3.0cm}p{7.5cm}}
 \caption{核心方法的来源与判断状态。}\label{tab:formula-origin}\\
 \toprule
 方法或公式 & 分类 & 当前解释 \\
 \midrule
 \endfirsthead
 \toprule
 方法或公式 & 分类 & 当前解释 \\
 \midrule
 \endhead
 有效工作/服务功能单位 & \literature & 已有计算系统和AI能效研究支持；本文扩展至制造业部署架构。 \\
 延迟容忍任务的时间迁移 & \literature、\adapted & 机制有数据中心灵活性文献；任务分类、截止窗和循环代表周是本研究设定。 \\
 PUE定义与设施能效 & \literature & 可依据ISO/IEC 30134-2；具体架构PUE仍需数据。 \\
 CPU/GPU成功任务吞吐路由 & \adapted & 以相同质量/SLA吞吐为原则；当前份额和倍数仍是结构情景。 \\
 任务数据与控制可接受性约束 & \adapted、\newwork & 数据驻留机制来自分布式计算的一般边界；具体任务--节点矩阵由本文定义，需由企业治理规则验证。 \\
 10\%或N+1取较高值的装机下限 & \newwork & 工程情景组合，不是已观测制造业平均规则。 \\
 三状态设施功率式 & \adapted & 状态划分有系统测量依据；当前线性组合与设施倍率是本研究实现。 \\
 安装IT容量与池化增益 & \newwork & 由已求解服务器存量和额定功率重构，用于解释集中化的统计复用，不是独立优化目标。 \\
 共用边缘能耗与架构选择能耗分离 & \newwork & 防止把所有架构共有的终端采集和控制能耗误算为部署收益；当前共用项尚待单独测量。 \\
 优化无AI净峰值基线 & \newwork & 用于阻止未观测余量或基线DER形成免费接入信用。 \\
 相对参考架构避免的接入容量 & \newwork & 仅报告MW和避免比例；不折算电网投资，不进入企业目标。 \\
 SIELS式批发边际价值 & \literature、\adapted & 只可作为独立后处理，不进入当前企业电费目标。 \\
 八项增量成本重构 & \implemented & 精确对应当前企业成本边界；不包含环境和电网升级项。 \\
 本地--云端联合任务分配与预留容量 & \implemented、\newwork & 已进入单行业结构敏感性；服务守恒和订阅容量约束对应代码，但条件触发的本地固定成本尚未实现。 \\
 运行用水与电网排放后处理 & \literature、\adapted & 因子乘法账户有文献基础；当前仅保留接口，地区与时序因子需另行校准。 \\
 \bottomrule
\end{longtable}

作者需优先判断：（1）是否将优化无AI净峰值基线和避免接入容量指标作为主要方法贡献；（2）装机下限是否保留10\%与N+1取较高值；（3）混合模型保留少量本地任务时，是否条件触发本地副本、N+1和固定运维成本；（4）联合优化中的云端产品是否继续采用预留CPU/GPU容量，还是另建不能与之混算的Token API模型；（5）批发边际价值是否进入补充敏感性；（6）未来区域节点模型是否仍属于同一篇论文。

\section{设施、能源与成本参数的估计方法}

本节说明各项物理和企业成本参数为什么需要、如何从来源转换到公式口径，以及哪些数值仍是情景。参数表的作用不是堆叠引用，而是让读者能够从服务器采购、功率状态、到户电价、DER技术和模型运维逐项重建成本边界。

\subsection{服务器采购、寿命、功率与设施倍率}

本地服务器基准为双NVIDIA L20 48~GB配置。采购价格从中国大陆公开服务器配置报价中选择与目标CPU、内存、存储和GPU数量最接近的规格；基准15万元/组，公开配置区间12.75--25.806万元/组\cite{lenovoL20}。高价配置同时改变CPU、内存和存储，因此该区间是整机采购情景，不是纯GPU价格区间。

服务器经济寿命基准5年，4--6年为主要敏感性。依据是阿里巴巴计算机设备和软件3--5年、青云服务器5年以及腾讯部分服务器会计寿命由4年调整至6年的公开披露。会计预计使用寿命不是物理寿命，但比任意折旧年限更接近企业采购口径。年度维护费取采购价5\%（2--8\%），软件、网络和小型机房附加投资取采购价20\%（10--35\%）；二者当前缺少匹配合同，均为D级作者假设。

双L20服务器满载、在线空闲和冷备墙插功率基准分别为1.50、0.36和0.02~kW/组；空闲与冷备区间分别为0.28--0.50和0.01--0.04~kW。L20单卡275~W规格来自HPE/NVIDIA资料，但GPU额定功率不能替代整机墙插测量\cite{hpeL20}。满载1.50~kW是包含CPU、内存、风扇和电源损耗的保守工程包络；三状态都必须用目标服务器在合格吞吐下实测替换。

本地边际设施倍率取1.30，低/高为1.20/1.60。雄安绿色信息化机房标准给出微型机房PUE不高于1.6、小型机房不高于1.3；LBNL小型服务器房研究观察到更高且高度离散的PUE\cite{xiongan2024,lbnlSmallRooms}。这里的1.30只表示新增AI负荷的冷却、UPS和其他辅助用电倍率，不是包含固定建筑负荷的全年机房PUE。

\begin{longtable}{p{3.5cm}p{2.3cm}p{2.4cm}p{4.2cm}p{1.2cm}}
 \caption{服务器与设施成本参数、范围和证据状态。}\label{tab:server-parameters}\\
 \toprule
 参数 & 基准 & 低--高 & 估计来源与用途 & 等级 \\
 \midrule
 \endfirsthead
 \toprule
 参数 & 基准 & 低--高 & 估计来源与用途 & 等级 \\
 \midrule
 \endhead
 双L20服务器采购价 & 15万元/组 & 12.75--25.806万元 & 联想公开整机配置；采购CAPEX & B \\
 经济寿命 & 5年 & 4--6年 & 中国互联网企业会计政策综合 & B \\
 年维护费 & 采购价5\% & 2--8\% & 维保合同缺失，情景参数 & D \\
 设施附加投资 & 采购价20\% & 10--35\% & 软件、网络和机房附加，情景参数 & D \\
 满载墙插功率 & 1.50~kW & 1.05--1.60~kW & L20规格加整机工程包络 & C \\
 在线空闲功率 & 0.36~kW & 0.28--0.50~kW & 部件测量与服务器包络综合 & C \\
 冷备待机功率 & 0.02~kW & 0.01--0.04~kW & BMC/软关机工程包络 & C \\
 边际设施倍率 & 1.30 & 1.20--1.60 & 雄安标准与LBNL小机房证据适配 & B/C \\
 贴现率 & 8\% & 6--12\% & 企业资本成本情景，非统计观测 & D \\
 \bottomrule
\end{longtable}

\subsection{企业电价与最大需量}

企业电价以广东珠三角五市1--10(20)~kV两部制工商业用户为地理和电压边界。2026年4月官方代理购电表给出平段0.65856875元/kWh、最大需量36.1元/kW/月，以及线损0.0147、输配0.1260、系统运行0.0859、政府性基金及附加0.02766875元/kWh\cite{gdTariff2026}。当前逐时情景用广东日前代表周电能量价格加上述体积附加；最大需量单列，不能平均进电量价格。平段官方电价仅作固定电价对照。

\boundary 该电价是广东特定电压等级和结算月份的企业付款代理，不代表全国制造企业合同。全国推广应按省份、用户类别、电压等级和市场交易方式更换参数；容量电价与最大需量计费通常二选一，不得同时相加。

\subsection{光伏、储能和屋顶容量}

屋顶光伏容量上限由

\begin{equation}
 \overline K_{i}^{\mathrm{PV}}=
 A_i^{\mathrm{roof}}u_i^{\mathrm{roof}}
 \eta^{\mathrm{module}}\gamma_i^{\mathrm{real}},
 \label{eq:pv-rooftop-capacity}
\end{equation}

估计，其中$A_i^{\mathrm{roof}}$为典型工厂屋顶面积，$u^{\mathrm{roof}}=0.90$为可用面积比例，$\eta^{\mathrm{module}}=0.22$~kW/m$^2$为组件功率密度，$\gamma^{\mathrm{real}}=0.80$为可实现比例。当前屋顶面积用美国EIA 2022 MECS同类行业的每企业封闭建筑面积作为中国敏感性代理\cite{mecs2022}，不是中国工厂屋顶测量，因此整体为C级并必须做跨国缩放敏感性。

当前有DER情景把光伏视为既有资产，增量CAPEX为零；这表示“已有屋顶光伏如何改变AI决策”，不表示新建光伏免费。若研究新建光伏，采用2.8元/W、寿命25年的情景成本并加入式~\eqref{eq:enterprise-with-der}，该造价仍需以目标年份中国工商业项目报价更新。

电池基准为2小时系统，往返效率0.96。能量部分99.189689欧元/kWh、逆变器60.296474欧元/kW，按7.8人民币/欧元换算；寿命分别22.5年和10年，逆变器固定运维为0.25\%/年。结构化成本表综合丹麦能源署技术目录和中国电力系统模板\cite{deaStorage}。由于其文字说明与结构化行对中国4小时系统总价的对应并不完全一致，当前值只能作为B/C级成本情景，正式结果需同时报告低/高成本与电池时长。

\subsection{模型状态与生命周期成本}

模型状态以Qwen2.5-32B-Instruct INT4/AWQ为代表。官方模型权重用于校验显存和存储下限；每部署点200~GB模型存储表示生产、测试、回滚和三份副本的工程包络，而不是单个权重文件大小\cite{qwen25}. 双L20服务器自带7.68~TB存储，因此当前不产生增量磁盘CAPEX。

每个主要版本4个参考GPU小时、每年1次主要更新、每部署点0.25名AI/MLOps全职当量和30万元/FTE年全成本用于初始化与运维账户。\newwork 其中4个GPU小时、0.25~FTE和人员成本均为D级情景参数，不是制造企业平均观测；运维成本可能主导存储成本，至少应报告0.10/0.25/0.75~FTE敏感性，并避免把普通IT人员和模型运维重复计费。

\section{环境足迹的独立后处理}

本节保留运行用水和电力间接环境影响的通用计算接口。它们的作用是把已求得的IT用电、电网购电和设施位置与水耗或排放因子相乘，用于比较部署位置改变后的运行足迹。该账户是物理结果后处理，不进入企业目标函数，也不被解释为完整生命周期评价\cite{xiao2025}。

令$P^{a,\mathrm{IT}}_{nt}$为应用设施倍率之前的IT功率，$p^{a,\mathrm{grid}}_{nt}$为购电功率，$\mathrm{WUE}^{a}_{nt}$为现场水使用效率，$\omega^{\mathrm{grid}}_{nt}$和$\gamma^{\mathrm{grid}}_{nt}$分别为电网单位购电的用水和碳因子，$\chi_E=10^3$用于在MW和kWh口径间转换，则

\begin{align}
 W^{a,\mathrm{direct}}
 &=\sum_{n,t}\chi_E\mathrm{WUE}^{a}_{nt}
 P^{a,\mathrm{IT}}_{nt}\Delta t,
 \label{eq:direct-water}\\
 W^{a,\mathrm{indirect}}
 &=\sum_{n,t}\chi_Ep^{a,\mathrm{grid}}_{nt}
 \omega^{\mathrm{grid}}_{nt}\Delta t,
 \label{eq:indirect-water}\\
 E^{a}_{\mathrm{CO_2}}
 &=\sum_{n,t}\chi_Ep^{a,\mathrm{grid}}_{nt}
 \gamma^{\mathrm{grid}}_{nt}\Delta t.
 \label{eq:grid-carbon}
\end{align}

直接用水由IT能量和现场WUE决定；间接用水和碳排放由电网购电与地区、时间相关因子决定。水稀缺加权结果应用位置相关的稀缺因子另行计算，并与体积用水分开报告。

\boundary 当前后处理尚忽略服务器、建筑、网络、光伏和电池的完整上游与报废足迹，也未统一校准不同地区和冷却方式的WUE。土地和建设材料使用独立的空间复用与绿地建设账户，不与式~\eqref{eq:direct-water}--\eqref{eq:grid-carbon}重复相加。

\section{未纳入当前主模型的扩展}

旧公式稿还曾列出多区域节点选址、多年服务器存量与退役、CVaR风险调整、多年净现值、池化打平条件和最优节点数。这些概念可用于未来扩展，但当前没有对应的位置候选集、服务器年份存量、不确定性分布、长期价格路径或网络与治理成本参数。因此，本SI不保留这些扩展的定量公式，也不将它们写成已实现方法。若后续启用，应在具备输入数据、校准和验证后作为独立模型版本重新定义。

\section{数据来源、处理与参数审计要求}

这一节用于建立从原始数据到模型参数的可追溯链。同一公式可能同时使用官方统计、企业披露和研究情景；因此，除了给出数值，还需要说明数据的样本、年份、地域、单位、处理和与公式变量的映射。表~\ref{tab:data-status}汇总当前已具备的说明和仍需补齐的数据。正式数据表应提供来源、样本、年份、空间范围、许可证、原始单位、清洗规则、缺失处理、最终样本量以及与公式变量的映射，而不是列出程序文件。

\begin{longtable}{p{3.3cm}p{6.0cm}p{5.4cm}}
 \caption{当前数据说明状态与高优先级缺口。}\label{tab:data-status}\\
 \toprule
 数据模块 & 当前估计与来源 & 结论前必须处理的缺口 \\
 \midrule
 \endfirsthead
 \toprule
 数据模块 & 当前估计与来源 & 结论前必须处理的缺口 \\
 \midrule
 \endhead
 企业数、就业与AI采用 & 第五次经济普查行业、规模和地区表；A级 & 规模以下企业采用率和任务条件份额 \\
 任务需求 & 业务驱动量乘采用、覆盖、频率、调用与服务规模 & 企业任务日志和2030联合情景验证 \\
 工业负荷 & EWELD为主，钢铁和运输设备外部曲线，三行业原型回退 & 中国行业年度电量来源、季节性和跨工厂方差 \\
 任务柔性 & 按任务SLA构造刚性/日内/批处理结构先验 & 制造企业队列、截止期和未服务任务数据 \\
 CPU/GPU计算 & 相同质量和SLA下吞吐原则，文献包络 & 目标模型与双L20/CPU服务器实测 \\
 服务器与设施 & 公开报价、厂商规格、机房标准与小机房研究 & 墙插三状态功率和本地设施倍率实测 \\
 企业电价 & 广东2026官方到户价和需量价 & 省份、电压等级、合同与市场交易异质性 \\
 电网接入容量 & 由匹配无AI反事实与AI净负荷峰值之差计算 & 真实合同容量、配变/馈线裕量和离散升级档位 \\
 PV与储能 & MECS屋顶代理和技术成本目录 & 中国工厂屋顶测量、项目报价、退化和循环寿命 \\
 生命周期 & 开源模型规格加初始化和运维情景 & 制造企业真实版本频率、存储和FTE \\
 环境因子 & 运行用水、电网间接用水和碳排放的独立后处理接口 & 地区--时序WUE、电网水耗和排放因子，以及完整生命周期边界 \\
 \bottomrule
\end{longtable}

每个参数应保留唯一标识、引用、访问日期、价格年、原始口径、转换公式和证据等级。若A级观测与D级情景共同进入同一复合参数，结果表应展示关键D级参数的敏感性；未经外部证据支持的值不得与观测值混列。

\section{验证、可复现性与当前限制}

验证的目的是区分“代码按预期实现了公式”和“模型已证明现实结果”。前者通过服务守恒、功率重构、成本复算和反事实一致性检查；后者还需要真实任务、硬件、企业电价和接入边界的外部校准。本节因此把实现验证、复现信息和现实推断限制分开报告。

\subsection{最低验证要求}

方法实现至少需要通过以下检查：

\begin{enumerate}
  \item 各架构有效服务守恒，刚性任务不延迟，柔性批次在截止窗内完成；
  \item 混合情景中本地与云端服务之和等于总服务，云份额不超过上限，逐时云端CPU/GPU需求不超过预留容量；
  \item CPU/GPU路由份额守恒，逐时计算量不超过在线容量；
  \item 安装容量不低于规划下限和模型状态下限；
  \item 三状态功率分项之和等于设施功率；
  \item 无AI和AI情景采用匹配的价格、基础负荷和DER规则；
  \item 新增接入容量等于AI净峰值超过无AI净峰值的正部分；
  \item 八项成本分项与行业等效增量总成本完全复算一致，且其中不含电网升级成本；
  \item 代表节点服务量乘恢复倍数后等于行业服务总量。
\end{enumerate}

\subsection{代码与复现信息的适当位置}

正式公开时，代码版本、运行环境、求解器、随机种子、输入哈希和生成时间应集中放在复现清单，而不是穿插于方法叙述。内部活动实现位于项目的核心模型模块和情景运行模块；公开投稿版应将本地路径替换为带版本号和DOI的代码、数据存档链接。

\subsection{当前推断限制}

当前方法是168小时代表周、单节点功率平衡和连续接入容量筛查。它尚未内生模拟线路潮流、变压器和馈线离散档位、N-1、季节性、容量扩展反馈、批发价格反馈或项目建设时序，也不把新增接入容量折算为电网升级成本。企业到户成本也不是社会资源成本。全国联合异构结果、真实工业任务吞吐、制造企业合同电价和真实接入边界仍是结论形成前的高优先级缺口。

\section{补充材料的后续更新顺序}

后续更新应按以下顺序进行：

\begin{enumerate}
  \item 补齐表~\ref{tab:data-status}列出的高优先级数据缺口；
  \item 冻结基础情景与敏感性参数，记录哪些是观测、文献代理或研究假设；
  \item 完成全国31行业$\times$三架构联合CPU/GPU重跑和验证；
  \item 按“建议报告的结果结构”一节统一生成结果表，不复用旧口径数值；
  \item 最后加入环境后处理、区域网络或批发边际价值敏感性，并保持其与当前八项企业成本分离。
\end{enumerate}

\begin{thebibliography}{00}

\bibitem{iso30134}
ISO/IEC, ``Information technology---Data centres---Key performance indicators---Part 2: Power usage effectiveness (PUE),'' ISO/IEC 30134-2.

\bibitem{siels2026}
S. Murthy \emph{et al.}, ``Simulation of Industrial Electrical Load Shapes (SIELS),'' manuscript ADAPEN-D-26-00491, submitted to \emph{Advances in Applied Energy}, 2026.

\bibitem{rockafellar2000}
R. T. Rockafellar and S. Uryasev, ``Optimization of conditional value-at-risk,'' \emph{Journal of Risk}, vol. 2, no. 3, pp. 21--41, 2000.

\bibitem{nbs2023}
国家统计局，《中国经济普查年鉴2023》，第五次全国经济普查综合卷表2--22、2--23、4--9和4--10，2025。[在线]. Available: \url{https://www.stats.gov.cn/sj/pcsj/jjpc/5jp/indexch.htm}

\bibitem{ibm2025}
IBM Institute for Business Value, ``From AI projects to profits,'' June 2025. [Online]. Available: \url{https://newsroom.ibm.com/2025-06-10-IBM-Study-Businesses-View-AI-Agents-as-Essential-Not-Just-Experimental}

\bibitem{eweld2023}
EWELD, ``Electricity load profiles of industrial and commercial users in South China,'' Figshare, version 3, 2023, doi: 10.6084/m9.figshare.21893808.v3.

\bibitem{uciSteel}
UCI Machine Learning Repository, ``Steel Industry Energy Consumption,'' doi: 10.24432/C52G8C.

\bibitem{ffeLoad}
Forschungsstelle f\"ur Energiewirtschaft, ``Normalized industrial electrical load profiles Germany.'' [Online]. Available: \url{https://opendata.ffe.de/dataset/normalized-industrial-electrical-load-profiles-germany/}

\bibitem{mlperf}
MLCommons, ``MLPerf Inference benchmarks.'' [Online]. Available: \url{https://mlcommons.org/benchmarks/inference-datacenter/}

\bibitem{AWSrecommender}
Amazon Web Services, ``SageMaker Inference Recommender.'' [Online]. Available: \url{https://docs.aws.amazon.com/sagemaker/latest/dg/inference-recommender.html}

\bibitem{lenovoL20}
联想，《ThinkSystem WR5220 G3 双NVIDIA L20服务器公开配置与报价》，访问于2026年8月。[在线]. Available: \url{https://m.lenovo.com.cn/wiki/product-doc-36784.html}

\bibitem{hpeL20}
Hewlett Packard Enterprise, ``NVIDIA Accelerators for HPE: QuickSpecs,'' rev. 74. [Online]. Available: \url{https://www.hpe.com/es/es/collaterals/collateral.c04123180.html}

\bibitem{xiongan2024}
雄安新区，《绿色信息化机房建设标准》，2024。[在线]. Available: \url{https://www.xiongan.gov.cn/20241223/e91521b5c735491b94c5c3e9d53b24c0/20241223e91521b5c735491b94c5c3e9d53b24c0_5178973c950dde4ce3bae27147fdf35274.pdf}

\bibitem{lbnlSmallRooms}
Lawrence Berkeley National Laboratory, ``Energy Efficiency in Small Server Rooms,'' 2013. [Online]. Available: \url{https://datacenters.lbl.gov/sites/default/files/SmallServerRooms_Final%20Report_0.pdf}

\bibitem{gdTariff2026}
广东电网，《2026年4月代理购电工商业用户价格表》，2026。[在线]. Available: \url{https://www.hengqin.gov.cn/attachment/0/432/432110/3896922.pdf}

\bibitem{mecs2022}
U.S. Energy Information Administration, ``2022 Manufacturing Energy Consumption Survey,'' Tables 7.6, 7.9 and 9.1. [Online]. Available: \url{https://www.eia.gov/consumption/manufacturing/data/2022/}

\bibitem{deaStorage}
Danish Energy Agency, ``Technology Data for Energy Storage,'' technology data catalogue; combined with the China power-system storage-cost template, 2025.

\bibitem{qwen25}
Qwen Team, ``Qwen2.5-32B-Instruct-AWQ model card,'' 2024. [Online]. Available: \url{https://huggingface.co/Qwen/Qwen2.5-32B-Instruct-AWQ}

\bibitem{xiao2025}
T. Xiao, F. Fuso Nerini, H. D. Matthews, M. Tavoni, and F. You, ``Environmental impact and net-zero pathways for sustainable artificial intelligence servers in the USA,'' \emph{Nature Sustainability}, vol. 8, pp. 1541--1553, 2025, doi: 10.1038/s41893-025-01681-y.

\end{thebibliography}

\section*{文档状态}

本文件是中文方法优先的SI活动底稿，不是投稿终稿。方法公式、参数估计路径和当前来源边界已按主流程统一；企业任务日志、目标硬件实测、真实接入容量边界、全国联合异构结果和外部验证仍待补充。任何后续数值更新都必须同时检查功能单位、无AI反事实、八项成本重构、新增接入容量指标、参数证据等级、结果表和正文结论是否一致。

\end{document}
